%!TEX root = paper.tex

\section{Module Extraction via Datalog Reasoning} \label{sec:module-settings}

 
In this section, we present our approach to module extraction
by reduction into a reasoning problem in datalog. 
Our approach builds on recent techniques
that exploit datalog engines for 
ontology reasoning 
\cite{DBLP:conf/ijcai/KontchakovLTWZ11,stefanoni2013introducing,DBLP:conf/aaai/ZhouNGH14}.
In what follows, we fix an arbitrary
TBox $\T$ and signature $\S\subseteq\sig{\T}$. Unless otherwise stated,
our definitions and theorems are parameterised
by such $\T$ and $\S$. We assume
w.l.o.g.\ that rules in $\T$ do not
share existentially quantified variables.
For simplicity, we also 
assume that $\T$ contains no constants 
(all our results can be
seamlessly extended). 

%%
\subsection{Overview and Main Intuitions}\label{sec:overview}
%%

Our overall strategy to
extract a module $\M$  of $\T$ 
for an inseparability relation $\equiv^z_{\S}$,
with $z \in \set{\mathsf{m},\mathsf{q},\mathsf{f},\mathsf{i}}$, 
can be summarised by the following steps:
% 
\begin{enumerate}
  \item Pick a substitution $\theta$ mapping all existentially 
  quantified variables in $\T$ to 
  constants, and transform $\T$ into a datalog program $\prog$ by 
  \begin{inparaenum}[\it(i)] 
  \item Skolemising 
all rules in $\T$ using $\theta$ and
\item turning disjunctions into conjunctions while splitting them into different rules,
thus replacing each function-free disjunctive rule of the form 
$\fml(\x) \rightarrow \bigvee_{i=1}^n\fmm_i(\x)$
with datalog rules 
$\fml(\x) \rightarrow \fmm_1(\x), \dots, \fml(\x) \rightarrow \fmm_n(\x)$.
  \end{inparaenum}

 \item Pick a $\Sigma$-ABox 
$\abox_0$ and materialise
$\prog \cup \abox_0$.
 \item Pick a set $\arel$ of ``relevant facts'' in the materialisation
 and compute the supporting rules in $\prog$ for
 each such fact.
\item 
The module $\M$ consists of all rules in $\T$ that
yield some supporting rule in $\prog$. In this way,
$\M$ is fully determined
by the substitution $\theta$ and the ABoxes $\abox_0$ and $\arel$.
\end{enumerate}
%
The main intuition behind our module extraction approach is that we
can pick $\theta$, $\abox_0$ and $\arel$ (and hence $\M$) such 
that each proof $\rho$ of a $\S$-consequence $\varphi$ of $\T$ to
be preserved can be embedded in a forward chaining proof $\rho'$
in $\prog \cup \abox_0$
of a relevant fact in $\arel$.
Such an embedding satisfies the key property that, for each rule $r$
involved in $\rho$, at least one corresponding datalog rule in $\prog$
is involved in $\rho'$. In this way we ensure that $\M$ contains the
necessary rules to entail $\varphi$.
%
This approach, however, does not
ensure minimality of $\M$: since 
$\prog$ is a strengthening of $\T$ there may be
proofs of a relevant fact in $\prog \cup \abox_0$ that do not
correspond to a $\S$-consequence of $\T$, which may lead to
unnecessary rules in~$\M$.

\begin{figure}
\[
{\small
\begin{array}{rl|r@{\:\sby\:}l}
(\ax_1) 	& \A(x) \rightarrow \exists y_1. [\roleR(x,y_1) \wedge \B(y_1)] 		& \A&\exists \roleR. \B \\
(\ax_2) 	& \A(x) \rightarrow \exists y_2. [\roleR(x,y_2) \wedge \C(y_2)] 		& \A&\exists \roleR. \C \\
(\ax_3) 	& \B(x)\wedge\C(x) \rightarrow \D(x) 											& \B\sqcap\C&\D \\
(\ax_4) 	& \D(x) \rightarrow \exists y_3. [\roleS(x,y_3) \wedge \E(y_3)] 		& \D&\exists \roleS. \E \\
(\ax_5) 	& \D(x) \wedge \roleS(x,y) \rightarrow \F(y) 									& \D&\forall \roleS. \F \\
(\ax_6) 	& \roleS(x,y) \wedge \E(y) \wedge \F(y) \rightarrow \G(x) 				& \exists \roleS. (\E\sqcap\F)&\G \\
(\ax_7) 	& \G(x) \wedge \H(x) \rightarrow \bot		& \G \sqcap \H& \bot 
\end{array}}\]
\caption{Example TBox $\Tex$ with DL translation}
\label{fig:exampleTBox}
\end{figure}

\begin{figure*}[t]
%  \centering
%  \begin{minipage}[t]{8cm}\vspace{0pt}
    \begin{small}
      \begin{tikzpicture}[parent anchor=south,sibling distance=15mm,level distance=10mm,>=stealth]
        \tikzstyle{level 2}=[sibling distance=18mm]
        \tikzstyle{level 3}=[sibling distance=15mm]
        \hskip20mm
        \node (top) {}
        child[grow=180,level distance=40mm] {node (leftTree) {$\G(a)$}
          [grow=-90,level distance=10mm]
          child {node (Safa1) {$\roleS(a,f(a))$}
            child {node {$\D(a)$}
              edge from parent node[right=-0.3mm] {$\ax'_4$}}}
          child[grow=-115,level distance=11mm]
          {node (Efa) {$\hskip5pt\E(f(a))$}
            child[grow=-90,level distance=10mm] {node {$\D(a)$}
              edge from parent node[right=-0.3mm] {$\ax''_4$}}
            edge from parent node[right=0mm] {$\ax_6$}}
          child {node (Ffa) {$\F(f(a))$}
            child {node (Safa2) {$\roleS(a,f(a))$}
              child {node {$\D(a)$}
                edge from parent node[right=-0.3mm] {$\ax'_4$}}}
            child {node {$\D(a)$}}}
          edge from parent[draw=none]}
        child[grow=0,level distance=40mm] {node (rightTree) {$\G(a)$}
          [grow=-90,level distance=10mm]
          child {node (Sac1) {$\roleS(a,c)$}
            child {node {$\D(a)$}}}
          child[grow=-115,level distance=11mm] {node (Ec) {$\hskip0pt\E(c)$}
            child[grow=-90,level distance=10mm] {node {$\D(a)$}}}
          child {node (Fc) {$\F(c)$}
            child {node (Sac2) {$\roleS(a,c)$}
              child {node {$\D(a)$}}}
            child {node {$\D(a)$}}}
          edge from parent[draw=none]};
        \node[below of=Ffa,node distance=5.7mm] {$\ax_5$};
        %\node[above right of=Safa1,node distance=2mm] (Safa1c) {};
        %\node[right of=Sac1,node distance=3mm] (Sac1c) {};
        \node[below of=leftTree,node distance=35mm] {(a)};
        \node[below of=rightTree,node distance=35mm] {(b)};
        % \hskip80mm
        % \node (right) {$\G(a)$}
        % child {node (Sac1) {$\roleS(a,c)$}
        %   child {node {$\D(a)$}}}
        % child[grow=-115,level distance=11mm] {node (Ec) {$\hskip0pt\E(c)$}
        %   child[grow=-90,level distance=10mm] {node {$\D(a)$}}}
        % child {node (Fc) {$\F(c)$}
        %   child {node (Sac2) {$\roleS(a,c)$}
        %     child {node {$\D(a)$}}}
        %   child {node {$\D(a)$}}};
        \draw[->,dotted,shorten <=-5pt,shorten >=-15pt] (Safa1.25) .. controls +(3,1) and +(-3,1) .. (Sac1.140) node[midway,below] {$\theta$};
%        \draw[->,dotted,shorten <=-5pt,shorten >=-14pt] (Ffa.35) .. controls +(3,1) and +(-3,1) .. (Fc.145);
        \draw[->,dotted,shorten <=-5pt,shorten >=-12pt] (Ffa.-35) .. controls +(3,-1) and +(-3,-1) .. (Fc.-140) node[near start,above] {$\theta$};
        \draw[->,dotted,shorten <=-5pt,shorten >=-12pt] (Safa2.-25) .. controls +(3,-1) and +(-3,-1) .. (Sac2.-125) node[midway,above] {$\theta$};
        \draw[->,dotted,shorten <=-5pt,shorten >=-12pt] (Efa.-35) .. controls +(3,-1) and +(-3,-1) .. (Ec.-140) node[midway,below] {$\theta$};
        \end{tikzpicture}
      \end{small}
%    \end{minipage}\hskip2mm%
  \caption{Proofs of $\G(a)$ from $\D(a)$ in (a) $\Tex$ and (b) the corresponding datalog program}
  \label{fig:proof-trees}
\end{figure*}


%%% Local Variables: 
%%% mode: latex
%%% TeX-master: "paper"
%%% End: 



To illustrate how our strategy might work in practice, 
suppose that $\T$ is $\Tex$
in Fig.\ \ref{fig:exampleTBox}, 
$\S =\set{\B,\C,\D,\G}$, and that
we want a module $\M$  that
is $\S$-implication inseparable from $\Tex$. This is a
 simple case 
since $\varphi=\D(x)\to\G(x)$ is
the only non-trivial $\S$-implication
entailed
by $\Tex$; thus, for $\M$ to be a module
we only require that $\M \models \varphi$.

Proving $\Tex \models \varphi$ 
amounts to proving $\Tex \cup \{\D(a)\} \models \G(a)$ (with $a$ a fresh constant).
Figure~\ref{fig:proof-trees}(a) depicts a hyper-resolution tree 
$\rho$ showing how $\G(a)$ can be derived
from the clauses corresponding to $\ax_4$--$\ax_6$ and $\D(a)$, with
rule $\ax_4$ transformed into clauses
%
\begin{align*}
\ax_4' = \D(x) \to \roleS(x,f(x_3)) & &  \ax_4'' = \D(x) \to \E(f(x_3))
\end{align*}
%
Hence $\M = \{\ax_4$--$\ax_6\}$ is a $\S$-implication 
inseparable module of $\Tex$, and as 
$\G(a)$ cannot be derived
from any subset of $\{\ax_4$--$\ax_6\}$, $\M$ is also minimal.

In our approach, we pick
$\abox_0$ to contain the initial fact $\D(a)$, $\arel$ to contain
the fact to be proved $\G(a)$, and we make $\theta$
map variable $y_3$ in $\ax_4$
to a fresh constant $c$, in which case rule
$\ax_4$ corresponds to the following datalog rules in $\prog$:
%
\begin{align*}
\D(x) \to \roleS(x,c) && \D(x) \to \E(c)
\end{align*}
%
Figure~\ref{fig:proof-trees}(b) depicts a forward chaining
proof $\rho'$ of $\G(a)$ in 
$\prog \cup \{\D(a)\}$.
As shown in the figure, $\rho$ can be embedded in $\rho'$ via $\theta$ by
mapping functional terms over $f$ 
to the fresh constant $c$. In this way, the rules
involved in $\rho$ are mapped to the datalog rules involved in $\rho'$
via $\theta$. Consequently, we will extract the (minimal) module $\M = \{\ax_4$--$\ax_6\}$. 

%%%%
\subsection{The Notion of Module Setting}
%%%%

The substitution $\theta$ and the ABoxes $\abox_0$
and $\arel$, which determine the extracted module, can be chosen
in different ways to ensure the preservation of different kinds of $\Sigma$-consequences.
The following notion of
a \bla{} captures in a declarative way the main elements of
our approach.


\begin{definition}
\label{def:bla}
A \emph{\bla} for $\T$ and $\S$ is a tuple 
$\upchi = \langle \theta, \abox_0, \arel \rangle$ 
with
$\theta$ a substitution from existentially quantified
variables in $\T$ to constants,
$\abox_0$ a $\Sigma$-ABox,
$\arel$ a \hbox{$\sig{\T}$}-ABox, and s.t.\
no constant in $\upchi$ occurs in $\T$.

 
The \emph{program} of $\upchi$ is the smallest datalog program
$\prog^{\chi}$ containing, for each
$\ax=\fml(\x)\to\exists\y.[\bigvee_{i=1}^n\fmm_i(\x,\y)]$ in $\T$, the
rule $\fml \rightarrow \bot$ if $n = 0$ and all rules $\fml \to
\fct\theta$ for each $1 \leq i \leq n$ and each atom $\fct$ in
$\psi_i$.  The \emph{support} of $\upchi$ is the set of rules $r \in
\prog^{\chi}$ that support a fact from $\arel$ in $\prog^{\chi} \cup
\abox_0$.  The \emph{module} $\M^{\chi}$ of $\upchi$ is the set of
rules in $\T$ that have a corresponding datalog rule in the support of
$\upchi$.
\end{definition}

%%%%
\subsection{Modules for each Inseparability Relation}\label{sec:all-modules}
%%%%

We next consider each
inseparability relation  $\equiv^z_{\S}$, where
$z \in \set{\mathsf{m},\mathsf{q},\mathsf{f},\mathsf{i}}$, and
formulate a specific setting $\upchi_z$
which provably yields
a $\equiv^z_{\S}$-module of $\T$. 
%In 
%Section \ref{sec:module-containment} we show that some of these settings
%are \emph{optimal}, in the sense that no other ``reasonable'' choice of
%the parameters $\theta$, $\abox_0$, and $\arel$ leads to smaller $\equiv^z_{\S}$-modules.

%%%%
\subsubsection{Implication Inseparability}

%%%%%
 
The example in Section~\ref{sec:overview}
suggests a natural setting $\upchiimp=\langle \theta, \abox_0, \arel \rangle$ 
that guarantees implication inseparability. 
As in our example, we pick 
$\theta$ to be as ``general'' as possible by
Skolemising each existentially quantified variable to 
a fresh constant.
For $\A$ and $\B$ predicates of the same arity $n$,
proving that $\T$ entails a $\S$-implication 
$\varphi=\A(x_1, \ldots, x_n)\to\B(x_1, \ldots, x_n)$, 
amounts
to showing that $\T \cup \{\A(a_1, \ldots, a_n)\} \models \B(a_1, \ldots, a_n)$
for fresh constants $a_1, \ldots, a_n$. Thus, following the ideas of our example,
we 
initialise $\abox_0$ with 
a fact  $\A(c_{\A}^1, \ldots, c_{\A}^n)$ for each
$n$-ary predicate $\A \in \S$, and $\arel$ with
a fact 
$\B(c_{\A}^1, \ldots, c_{\A}^n)$
for each pair of $n$-ary predicates $\{\B,\A\} \subseteq \S$ with 
$\B \neq \A$.
% In this way we can ensure
% that all proofs of $\varphi$ in $\T$ can 
% be embedded in a proof of a relevant fact
% in $\prog^{\chiimp} \cup \{\A(\vec{c_{\A}})\} $, 
% and hence $\M^{\chiimp}$ entails the same
% $\S$-implications as $\T$.

\begin{definition} \label{def:chiimp} %
For each existentially quantified variable $y_j$ in $\T$, let $c_{y_j}$
be a fresh constant. Furthermore, 
for each $\A\in\S$ of arity $n$, let $c_{\A}^1, \ldots, c_{\A}^n$ be 
also fresh constants.  The setting
  $\upchiimp=\langle \theta^{\mathsf{i}}, \abox_0^{\mathsf{i}}, \arel^{\mathsf{i}} \rangle$  is defined as follows:
  \begin{itemize}
  \item $\theta^{\mathsf{i}}=\mset{y_j\mapsto c_{y_j}}{y_j \text{ existentially quantified in }
      \T}$,
  \item $\abox_0^{\mathsf{i}}=\{\A(c_{\A}^1, \ldots, c_{\A}^n) \mid \A \text{~$n$-ary predicate in } \S\}$, and
  \item $\arel^{\mathsf{i}}\,{=}\,\{\B(c_{\A}^1, \ldots, c_{\A}^n)\,{\mid}\,
    \A \neq \B\, \text{~$n$-ary predicates in } \S \}$.\qedhere 
  \end{itemize}  
\end{definition}
%
%\noindent
%
The setting
$\upchiimp$ is reminiscent of the datalog encodings 
typically used  to check whether a
concept  $A$ is subsumed by concept $B$
w.r.t.\  a  ``lightweight'' 
ontology $\T$ \cite{DBLP:conf/semweb/KrotzschRH08,stefanoni2013introducing}.
There, variables in rules are Skolemised as
fresh constants to produce a datalog program $\prog$
and it is then checked whether $\prog \cup \{A(a)\} \models B(a)$.
% The following theorem shows that, as expected, $\upchiimp$ 
% yields a $\eqi_{\S}$-module.
 
\begin{restatable}{theorem}{datalogAtomicSubsumptionModules}
\label{thm:datalogAtomicSubsumptionModules}
$\M^{\chiimp}\eqi_{\S} \T$.  
\end{restatable} 
 

%%%%
\subsubsection{Fact Inseparability}
%%%%

The setting $\upchiimp$ in
Def.~\ref{def:chiimp} cannot be used to ensure fact inseparability.
Consider again $\Tex$ and $\S = \set{\B,\C, \D,\G}$, for which
$\M^{\chiimp} = \{\ax_4, \ax_5, \ax_6\}$. For $\abox =
\set{\B(a),\C(a)}$ we have $\Tex \cup \abox \models \G(a)$ but
$\M^{\chiimp} \cup \abox \not\models \G(a)$, and hence $\M^{\chiimp}$
is not fact inseparable from $\Tex$.

More generally, 
$\M^{\chiimp}$ is only guaranteed to preserve
$\S$-fact entailments \mbox{$\T\cup\abox\models\fct$} where 
$\abox$ is a singleton. However, for a module to be fact inseparable
from $\T$ it must preserve all $\S$-facts when coupled with
\emph{any} $\S$-ABox.   
 We achieve this by choosing 
 $\abox_0$ to be the \emph{critical ABox}
for~$\S$, which consists of all facts that can be
constructed using $\S$ and a single fresh constant
\cite{DBLP:conf/pods/Marnette09}. Every $\S$-ABox
can be homomorphically mapped into the critical
$\S$-ABox. 
In this way, we can show
that all proofs of a $\S$-fact in $\T \cup \abox$ can 
be embedded in a proof of a relevant fact in $\prog^{\chi} \cup \abox_0$.



\begin{definition} \label{def:chidd} %
  Let constants $c_{y_i}$ be as in Def.\ \ref{def:chiimp},
  and let $*$ be a fresh constant. The setting
  $\upchidd=\langle \theta^{\mathsf{f}}, \abox_0^{\mathsf{f}}, \arel^{\mathsf{f}} \rangle$ is defined as follows:
  \begin{inparaenum}[\it (i)]
  \item $\theta^{\mathsf{f}} = \theta^{\mathsf{i}}$,
  % \item $\abox_0^{\mathsf{f}}$ consists of all facts $\A(\ast, \ldots, \ast)$
  % with $\A$ an $n$-ary $\S$-predicate, and
  \item $\abox_0^{\mathsf{f}} = \mset{\A(\ast,\ldots,\ast)}{\A\in\S}$, and
  \item $\arel^{\mathsf{f}} = \abox_0^{\mathsf{f}}$\qedhere
  \end{inparaenum}
\end{definition}
%
The datalog programs for $\upchiimp$ and  $\upchidd$ coincide
and hence the only difference between the two settings is in the definition of
%the initial ABox.
their corresponding ABoxes.  In our example, both
$\abox_0^{\mathsf{f}}$ and $\arel^{\mathsf{f}}$ contain facts
$\B(\ast)$, $\C(\ast)$, $\D(\ast)$, and $\G(\ast)$. Clearly,
$\prog^{\chidd} \cup \abox_0 \models \G(\ast)$ and the 
proof additionally involves rule $\ax_3$. Thus $\M^{\chidd} = \set{\ax_3,\ax_4,\ax_5,\ax_6}$.
%  The following theorem establishes that
% $\M^{\chidd}$ is indeed a $\eqf_{\S}$-module of $\T$.

\begin{restatable}{theorem}{disjDatalogModules}
\label{thm:disjDatalogModules}
$\M^{\chidd}\eqf_{\S}\T$.
\end{restatable}
%

%%%
\subsubsection{Query Inseparability}
%%%
Positive existential queries constitute a much richer query 
language than facts as they
allow for existentially quantified variables. 
Thus, the
query inseparability requirement  invariably 
leads to larger modules.

For instance, let $\T=\Tex$ and
$\S=\set{\A,\B}$. Given the $\S$-ABox
$\abox = \{\A(a)\}$ and $\S$-query $q = \exists y.\B(y)$ we have that
$\Tex \cup \abox \models q$ (due to rule $\ax_1$). 
The module
$\M^{\chidd}$ is, however, empty. Indeed, the materialisation 
of $\prog^{\chidd} \cup \{\A(\ast)\}$ consists of the additional facts
$\roleR(\ast,c_{y_1})$ and $\B(c_{y_1})$ and hence it does not
contain any relevant fact mentioning only $\ast$.
Thus, 
$\M^{\chidd} \cup \abox \not\models q$ and $\M^{\chidd}$ is not
query inseparable from $\Tex$.

Our example suggests that, although the critical ABox
is constrained enough to embed every $\S$-ABox,  we may need
to consider additional relevant facts to capture all
proofs of a $\S$-query. 
In particular, rule $\ax_1$ implies that $\B$ contains an 
instance
whenever $\A$ does: a dependency that is then checked by $q$.
This can be captured by considering fact $\B(c_{y_1})$
as relevant, in which case rule $\ax_1$ would be in the module.

More generally, we consider a module setting $\upchi$ 
that differs from $\upchidd$
only in that all $\S$-facts (and not just those over 
$\ast$) are considered
as relevant. 
%
\begin{definition} \label{def:chicq} %
  Let constants $c_{y_i}$ and $\ast$ be as in Def.~\ref{def:chidd}. 
  The setting $\upchicq=\langle \theta^{\mathsf{q}}, \abox_0^{\mathsf{q}}, \arel^{\mathsf{q}} \rangle$ is 
  as follows:
  \begin{inparaenum}[\it (i)]
  \item $\theta^{\mathsf{q}}= \theta^{\mathsf{f}}$,
  \item $\abox_0^{\mathsf{q}} = \abox_0^{\mathsf{f}}$, and 
  \item $\arel^{\mathsf{q}}$ consists of all $\S$-facts $\A(a_1, \ldots, a_n)$
  with each $a_j$ either a constant $c_{y_i}$ or $\ast$.\qedhere
  \end{inparaenum}
\end{definition}
%
Correctness is established by the following theorem:


\begin{restatable}{theorem}{datalogCQModules} \label{thm:datalogCQModules}
  $\M^{\chicq}\eqq_{\S}\T$.
\end{restatable}
 
%% 
\subsubsection{Model Inseparability}
\label{sec:semantic-modules}

The modules generated by $\chicq$ may not be model inseparable from
$\T$. To see this, let $\T=\Tex$ and $\S = \{\A,\D,\roleR\}$, in which case
$\M^{\chicq}=\set{\ax_1,\ax_2}$. The interpretation 
$\calI$ where 
$\Delta^{\calI} = \{a,b\}$, $\A^{\calI} = \{a\}$,
$\B^{\calI} = \C^{\calI} = \{b\}$, $\D^{\calI} = \emptyset$
and $R^{\calI} = \{ (a,b)\}$ is a model of
$\M^{\chicq}$. This interpretation, however, 
cannot be extended to a model of $\ax_3$ (and hence of $\T$) without reinterpreting
$\A$, $\roleR$ or $\D$. 

%
The main insight behind locality and reachability
modules
%\cite{GrauHKS08JAIR,DBLP:conf/lpar/NortjeBM13},
is to ensure that each model of the module can be extended
to a model of $\T$ in a uniform way. 
Specifically,
each model of a $\top\!\bot\!^*$-locality or
$\top\!\bot\!^*$-reachability module can be
extended to a model of $\T$ by interpreting all other
predicates $\A$ as either $\emptyset$ or
$(\Delta^{\calI})^{n}$ with $n$ the arity of $\A$.
Thus, $\M = \{\ax_1,\ax_2,\ax_3\}$ is
a $\eqm_{\S}$-module of $\Tex$ since all its models
can be extended
by interpreting $\E$, $\F$ and $\G$ as
the domain, $\H$ as empty, and $\roleS$ as the Cartesian 
product of the domain.
We can capture this idea in our framework by means of the following setting.
%
\begin{definition} \label{def:chistar}
  The setting $\upchistar=\langle \theta^{\mathsf{m}}, \abox_0^{\mathsf{m}}, \arel^{\mathsf{m}} \rangle$ is
  as follows:
  $\theta^{\mathsf{m}}$ maps  each existentially quantified variable  to the fresh constant
  $\ast$ and
  $\abox_0^\mathsf{m} = \arel^\mathsf{m} = \abox_0^{\mathsf{f}}$. %\mset{\A(\ast, \ldots, \ast)}{\A\in\S}$.
\end{definition}
%
In our example, 
$\prog^{\chistar}\cup\abox_0^{\mathsf{m}}$ entails the
relevant facts
$\A(\ast)$, $\roleR(\ast,\ast)$ and $\D(\ast)$, and hence  
$\M^{\chistar} = \{\ax_1,\ax_2,\ax_3\}$.

To show that $\M^{\chistar}$
is a $\eqm_{\S}$-module we prove that
all models $\calI$ of $\M^{\chistar}$
can be extended to a model of $\T$
as follows:
% 
\begin{inparaenum}[\it(i)]
\item predicates not occurring in the materialisation
of $\prog^{\chistar} \cup \abox_0^{\mathsf{m}}$ are interpreted as empty;
\item predicates in the support
of $\upchistar$ (and hence occurring in $\M^{\chistar}$)
are interpreted as in $\calI$; and
\item all other predicates $\A$ are interpreted as $(\Delta^{\calI})^{n}$ with $n$ the arity of $\A$. 
\end{inparaenum}
%Thus, we obtain the following theorem.
%
\begin{restatable}{theorem}{datalogSemModules} \label{thm:datalogSemModules}
  $\M^{\chistar}\eqm_{\S}\T$. 
\end{restatable}



%%%%%
\subsection{Modules for Ontology Classification}\label{sec:modules-class}

Module extraction has been 
exploited for 
optimising ontology classification~\cite{DBLP:conf/semweb/RomeroGH12,DBLP:conf/ore/TsarkovP12,DBLP:journals/jar/GrauHKS10}.
In this case, it is not only required that 
modules are implication inseparable from $\T$, but also that they
preserve all implications  $\A(\x) \to \B(\x)$ with 
$\A\in \S$ but $\B \notin \S$. This requirement
can be captured as given next.

\begin{definition}
  TBoxes $\T$ and $\T'$ are \emph{$\S$-classification inseparable}
  ($\T\eqc_{\S}\T'$) if for each $\fml$ of the form $\A(\x)\to\B(\x)$
  %or $\A(\x)\to$ 
  with $\A\in\S, $ and
  $\B\in\sig{\T\cup\T'}$ we have $\T\models\fml$ iff $\T'\models\fml$.
 \end{definition}
% 
Classification inseparability is a stronger requirement
than implication inseparability. For
$\T = \{ \A(x) \to \B(x)\}$ and $\S = \{ \A \}$,
$\M = \emptyset$ is implication inseparable from $\T$, whereas
classification inseparability requires that $\M = \T$. 

Modular reasoners such as MORe and Chainsaw rely on
locality $\bot$-modules, which satisfy this requirement.
Each model of a $\bot$-module $\M$ can be extended
to a model of $\T$ by interpreting all additional predicates as empty, which
is not possible if $\A \in \S$ and $\T$ entails $\A(x) \to \B(x)$ but $\M$ does not.
We can cast $\bot$-modules in our framework with
the following setting, which extends $\upchistar$
in Def.~\ref{def:chistar} by also considering
as relevant facts involving predicates not in $\S$.
%
\begin{definition} \label{def:chibot}
  The setting $\upchibot=\langle \theta^{\mathsf{b}}, \abox_0^{\mathsf{b}}, \arel^{\mathsf{b}} \rangle$ is as
  follows: $\theta^{\mathsf{b}} = \theta^{\mathsf{m}}$, 
  $\abox_0^{\mathsf{b}} = \abox_0^{\mathsf{m}}$, and 
  $\arel$ consists of all facts $\A(\ast, \ldots, \ast)$ where
  $\A\in\sig{\T}$.\qedhere
\end{definition}
%
The use of $\bot$-modules is, however, 
stricter than is needed for ontology classification.
For instance, if we consider
$\T = \Tex$ and $\S = \{\A\}$ we have
that $\M^{\chibot}$ contains all rules
$\ax_1$--$\ax_6$, but since $\A$ does not
have any subsumers in $\Tex$ the empty TBox
is already classification inseparable from $\Tex$.

The following module setting extends
$\upchiimp$ in Def.\ \ref{def:chiimp}
to ensure classification inseparability.
As in the case of $\chibot$ in Def.~\ref{def:chibot}
the only required modification is to
also consider as relevant facts 
involving predicates outside $\S$.

 
\begin{definition} \label{def:chicls}
 Setting $\upchicls= (\theta^{\mathsf{c}}, \abox_0^{\mathsf{c}}, \arel^{\mathsf{c}})$
  is as follows:
  $\theta^{\mathsf{c}} = \theta^{\mathsf{i}}$,
  $\abox_0^{\mathsf{c}} = \abox_0^{\mathsf{i}}$,
  and $\arel^{\mathsf{c}}$ consists of all facts
  $\B(c_{\A}^1, \ldots, c_{\A}^n)$ s.t.\ 
  $\A \neq \B$ are $n$-ary predicates, $A \in \S$ and
  $\B\in\sig{\T}$.
\end{definition}
%
Indeed, if we consider again $\T = \Tex$ and $\S = \{\A\}$, the module
for $\upchicls$ is empty, as desired.

\begin{restatable}{theorem}{datalogClassifModules}
  \label{thm:datalogClassifModules}
  $\M^{\chicls}\eqc_\S\T$. 
\end{restatable}

%%%%
\subsection{Additional Properties of Modules}
%%%%

Although the essential property of a module $\M$ is
that it captures all relevant $\S$-consequences of $\T$,
in some applications it is
desirable that modules satisfy additional requirements.
  
In ontology reuse scenarios, it is sometimes desirable that a module
$\M$ does not ``leave any relevant information behind'', in the sense that
$\T \setminus \M$ does not
entail any relevant $\S$-consequence---a property 
referred to as \emph{depletingness}
\cite{DBLP:journals/ai/KontchakovWZ10}.
%
\begin{definition}
  Let $\equiv_{\S}^z$ be an inseparability relation.
  A $\equiv_{\S}^z$-module $\M$ of $\T$ is \emph{depleting} if
  $\T\setminus\M\equiv_{\S}^z\emptyset$.
  %\footnote{In the literature, a stronger notion 
  %is often considered where it is required that $\T\setminus\M\equiv_{\S \cup \sig{\M}}^z\emptyset$.
  %The motivation for this strengthening is preservation of justifications. Although our modules are not
  %strongly depleting, they do preserve justifications.}
\end{definition}
%
Note that not all modules are depleting: for some relevant
$\S$-entailment $\varphi$ it may be that $\M \models \varphi$
(as required by the definition of module), but also that
$(\T \setminus \M) \models \varphi$, in which case $\M$ is not depleting.
The following theorem establishes 
that all modules defined in Section \ref{sec:all-modules}
are depleting.
%
\begin{restatable}{theorem}{depletingness}
  $\M^{\chi_z}$ is depleting for each $z\in\set{\mathsf{m},\mathsf{q},\mathsf{f},\mathsf{i},\mathsf{c}}$. 
\end{restatable}


Another common application of modules is to optimise the
computation of justifications: minimal subsets of
a TBox that are sufficient to entail a given formula
\cite{DBLP:conf/semweb/KalyanpurPHS07,DBLP:conf/aswc/SuntisrivarapornQJH08}.
%
\begin{definition}
  Let $\T\models\fml$. A \emph{justification
    for $\fml$ in $\T$} is a minimal subset $\T'\subseteq\T$ such that
  $\T'\models\fml$.
\end{definition}
%
Justifications are displayed in ontology development
platforms as explanations of why an entailment holds, and 
tools typically compute all of them. Extracting justifications
is a computationally intensive task, and locality-based modules have
been used to reduce the size of the problem: if $\T'$
is a justification of $\varphi$ in $\T$, then $\T'$ is contained in
a locality module of $\T$ for  $\S = \sig{\varphi}$. 
Our modules are
also justification-preserving, and we can adjust
our modules depending on what kind of first-order sentence $\varphi$ is.
%
\begin{restatable}{theorem}{presJustifications} \label{thm:presJustifications} ~
 Let $\T'$ be a justification for a first-order sentence $\varphi$ in $\T$ and let
 $\sig{\varphi} \subseteq \S$. Then, $\T' \subseteq \M^{\chistar}$. Additionally,
 the following properties hold:
 %
 \begin{inparaenum}[\it (i)]
 \item if $\varphi$ is a rule, then $\T' \subseteq \M^{\chicq}$;
 \item if $\varphi$ is datalog, then $\T' \subseteq \M^{\chidd}$; and
 \item if $\varphi$ %is an implication,
   is of the form $\A(\x)\to\B(\x)$, then $\T' \subseteq
   \M^{\chiimp}$; finally, if $\varphi$ satisfies
   $A \in \S, B \in \sig{\T}$, then
   $\T' \subseteq \M^{\chicls}$.
 \end{inparaenum}
\end{restatable}



%%%
\subsection{Complexity of Module Extraction}
\label{sec:complexity}
%%%
 
We conclude this section by showing that our modules can be 
efficiently computed in most practically relevant cases. 

\begin{restatable}{theorem}{complexity} \label{thm:complexity} %
  Let $m$ be a non-negative integer and $L$ a class of TBoxes s.t.\ 
  each rule in a TBox from $L$ has at most $m$
  distinct universally quantified variables.  The following problem is
  tractable: given $z \in \{\mathsf{q}, \mathsf{f}, \mathsf{i}, \mathsf{c}\}$, $\T\in L$, and 
  $\ax\in\T$, decide whether \mbox{$\ax\in\M^{\chi_z}$}.
 The problem is solvable in polynomial time for arbitrary classes $L$ of TBoxes if
 $z = \mathsf{m}$. 
\end{restatable}

We now provide a proof sketch for this result.
Checking whether a datalog program $\prog$ and an ABox $\abox$
entail a fact
is feasible in  $\mathcal{O}(\vert \prog \vert \cdot n^v)$, with 
$n$ the number 
of constants in $\prog \cup \abox$ and $v$ the maximum
number of variables in a rule from $\prog$ \cite{DBLP:journals/csur/DantsinEGV01}.
Thus, although datalog reasoning is exponential in the size of $v$
(and hence of $\prog$), it
is tractable if $v$ is bounded by a constant.

Given arbitrary $\T$ and $\S$, and for $z \in \{\mathsf{m}, \mathsf{q}, \mathsf{f}, \mathsf{i}, \mathsf{c}\}$,
the datalog program $\prog^{\chi_z}$ 
can be computed in linear time in the size of $\vert \T \vert$.
The number of constants $n$ in $\upchi_z$ (and hence in
$\prog^{\chi_z} \cup \abox_0^z$) is
linearly bounded in $\vert \T \vert$, whereas the maximum number of 
variables $v$ coincides with the maximum number
of universally quantified variables in a rule from $\T$. 
As shown in \cite{DBLP:conf/aaai/ZhouNGH14}, 
computing the support of a fact in a datalog program is
no harder than fact entailment, and thus module extraction in our approach is
feasible in  $\mathcal{O}(\vert \T \vert \cdot n^v)$, and thus
tractable for ontology languages where rules have a bounded number of variables (as is the case for
most DLs). Finally, if $z = \mathsf{m}$ the setting $\upchistar$ involves 
a single constant $\ast$ and module extraction boils down to reasoning
in propositional datalog (a tractable problem regardless of $\T$).





%%%
\subsection{Module Containment and Optimality} \label{sec:module-containment}
%%

Intuitively, the more expressive the language for which preservation
of consequences is required the larger modules need to be.
The following proposition shows 
that our modules are consistent with this 
intuition.
% For this, we use the following notion of a homomorphism between
% module settings.
%
\begin{restatable}{proposition}{hierarchy} \label{prop:module-hierarchy}
  $\M^{\chiimp} \subseteq \M^{\chidd} \subseteq \M^{\chicq} \subseteq
  \M^{\chistar}\subseteq\M^{\chibot}$ and $\M^{\chiimp} \subseteq \M^{\chicls} \subseteq \M^{\chibot}$
\end{restatable}
%
As already discussed, these containment relations 
are strict for many $\T$ and $\S$.

We conclude this section by discussing whether 
each
$\upchi_z$ with $z \in \{\mathsf{q}, \mathsf{f}, \mathsf{i}, \mathsf{c}\}$
is optimal for its inseparability relation in the sense that there is no
setting that produces smaller modules.
To make optimality statements precise we need to 
consider families of module settings, that is, functions
that assign a module setting for each pair of $\T$ and $\S$.
%
\begin{definition}
  A \emph{setting family} is a function $\Psi$ that maps a TBox
  $\T$ and signature $\S$ to a \bla{} for $\T$ and $\S$. 
  We say that $\Psi$ is \emph{uniform} if for every $\S$ and
  pair of TBoxes $\T,\T'$ with the same number of
  existentially
  quantified variables $\Psi(\T,\S)=\Psi(\T',\S)$.
  Let $z\in\set{\mathsf{i},\mathsf{f},\mathsf{q},\mathsf{c}}$; then,
  $\Psi$ is \emph{$z$-admissible} if, for each
  $\T$ and $\S$, $\M^{\Psi(\T,\S)}$ is a $\equiv^z_\S$-module of $\T$. 
  Finally, $\Psi$ is \emph{$z$-optimal} if
  $\M^{\Psi(\T,\S)}\subseteq\M^{\Psi'(\T,\S)}$ for every $\T$, $\S$
  and every uniform $\Psi'$ that is $z$-admissible.
\end{definition}
%
Uniformity ensures that settings
do not depend on the specific shape of rules in $\T$, but rather only on $\S$
and the number of existentially quantified variables in $\T$.
In turn, admissibility ensures that each setting yields a module.
The (uniform and admissible) family $\Psi^z$
corresponding to each setting  $\upchi^z$
in Sections \ref{sec:all-modules} and \ref{sec:modules-class}
is defined in the obvious way: 
for each $\T$ and $\S$, $\Psi^z(\T,\S)$ is the
setting $\upchi^z$ for $\T$ and $\S$.

The next theorem shows that $\Psi^z$ is optimal for implication and
classification inseparability.
%
\begin{restatable}{theorem}{optimality} \label{thm:optimality} %
  $\Psi^z$ is $z$-optimal for $z\in\set{\mathsf{i},\mathsf{c}}$.
\end{restatable}
%
In contrast, $\Psi^{\mathsf{q}}$ and  $\Psi^{\mathsf{f}}$
are not optimal. To see this, let
$\T = \{ \A(x) \to \B(x), \B(x) \to \A(x)\}$ and $\S = \{\A\}$. 
The empty TBox is fact inseparable from $\T$ since the only $\S$-consequence
of $\T$ is the tautology $\A(x) \to \A(x)$. However, $\M^{\chidd} = \T$ since 
fact $\A(a)$ is in $\arel^{\mathsf{f}}$ and its support is included in the module.
We can provide a family of settings that distinguishes
tautological from non-tautological inferences (see 
%extended version of the paper);
appendix);
%supplemental material);
% We can provide a family of settings that distinguishes
% tautological from non-tautological inferences (due to space restrictions we do
% not include any details about this family here);
however, this family yields settings of exponential size in $\vert \T \vert$, which is
undesirable in practice.

%%% Local Variables: 
%%% mode: latex
%%% TeX-master: "paper"
%%% End: 





%%% Local Variables: 
%%% mode: latex
%%% TeX-master: "paper"
%%% End: 
